\chapter{Pendahuluan}

%% VERSI BAHASA INDONESIA.
%%
%% KONVENSI PERISTILAHAN (berlaku untuk seluruh bab):
%%  1. Prosa akademik penghubung ditulis dalam bahasa Indonesia.
%%  2. Istilah teknis DIPERTAHANKAN dalam bahasa Inggris (ditulis miring) kecuali
%%     padanannya benar-benar mapan dan lazim dipakai. Menerjemahkan istilah yang
%%     belum mapan menghasilkan rumusan yang terasa janggal dan justru mengaburkan
%%     rujukan ke literatur -- misalnya "observation gate" dan "faithfulness"
%%     dipertahankan, bukan diterjemahkan.
%%  3. Yang diterjemahkan hanya istilah dengan padanan mapan: rekam medis elektronik,
%%     panduan klinis, basis data, presisi, keselarasan, kohort.
%%  4. Nama standar, produk, dan akronim tidak diterjemahkan
%%     (HL7 FHIR, SNOMED CT, RxNorm, LOINC, UMLS, Neo4j, Synthea).
%%  5. Akronim diperkenalkan sebagai: istilah Indonesia (Istilah Inggris, AKRONIM).
%%
%% \crosscheck{...} menandai isi yang wajib diverifikasi manusia sebelum diserahkan.

\section{Latar Belakang}

%% Catatan lingkup: beban penyakit pada skala populasi tidak dilaporkan pada tingkat CCD
%% -- CCD adalah kategori tata laksana klinis yang diperkenalkan \cite{virani2023}, bukan
%% kategori penyebab dalam Global Burden of Disease. IHD, yang fase kronisnya adalah CCD,
%% merupakan tingkat terdekat yang dilaporkan, sehingga angka di bawah ini bersifat
%% spesifik IHD, bukan kardiovaskular secara keseluruhan.
Penyakit jantung koroner kronis (\emph{chronic coronary disease}, CCD)---sebelumnya
dikenal sebagai penyakit jantung iskemik stabil (\emph{stable ischemic heart disease},
SIHD)---adalah fase kronis penyakit jantung iskemik (\emph{ischemic heart disease},
IHD), yang merupakan penyumbang terbesar untuk penyakit kardiovaskular global. \emph{Global
Burden of Disease Study 2019} memperkirakan terdapat 197 juta prevalensi IHD dan
9,14 juta kematian akibat IHD di dunia pada tahun 2019---dari total 18,6 juta kematian
kardiovaskular---dan beban tersebut terus meningkat di hampir seluruh negara yang bukan dalam
kelompok berpenghasilan tinggi \cite{gbd2019_cvd}. Indonesia mengikuti pola yang sama:
antara tahun 2000 dan 2019, angka kematian kardiovaskular terstandardisasi usia
(\emph{age-standardised mortality rate}) di Indonesia
naik dari 356,05 menjadi 412,46 kematian per 100.000 penduduk laki-laki,
sementara pada perempuan pada dasarnya tidak berubah \cite{gbd_indonesia_cvd}.
Pasien dengan CCD memerlukan pencegahan sekunder seumur hidup,
seperti terapi medis sesuai panduan (\emph{guideline-directed medical therapy}, GDMT) yang terbukti
menurunkan risiko kejadian kardiovaskular mayor (\emph{major adverse cardiovascular events},
MACE) secara bermakna. Untuk membakukan tata laksana tersebut, American Heart Association dan American College of
Cardiology menerbitkan panduan \emph{2023 AHA/ACC Guideline for the Management of Patients
with Chronic Coronary Disease} \cite{virani2023}. Panduan itu merangkum bukti ilmiah ke
dalam sekumpulan rekomendasi. Cakupannya meliputi terapi antiplatelet, pengelolaan lipid dan
tekanan darah, penyekat beta (\emph{beta blocker}), penghambatan
renin-angiotensin-aldosteron, pengelolaan berat badan dan gaya hidup, serta kontraindikasi
keamanan dari terapi-terapi tersebut.

Meskipun panduan klinis tersebut tersedia, penerapan dalam praktik rutin masih jauh dari optimal.
Studi INTERASPIRE, yang meneliti 4.548 pasien penyakit jantung koroner di 14 negara pada
keenam wilayah WHO, menemukan bahwa hanya 1,0\% pasien yang memenuhi definisi
penerapan panduan klinis yang optimal; hanya 38,6\% mencapai tekanan darah di bawah
130/80~mmHg, 16,6\% mencapai target kolesterol LDL di bawah 1,4~mmol/L, dan
9,0\% mengikuti rehabilitasi jantung \cite{adherence_gap}.
Klinisi harus memadukan rekam medis elektronik (\emph{electronic health record}, EHR) yang
terfragmentasi, puluhan kondisi dan obat, serta panduan klinis yang memuat banyak
rekomendasi bersyarat dan bersifat spesifik terhadap jenis pasien---seringkali di bawah tekanan waktu yang berat.
Sistem pendukung keputusan klinis (\emph{clinical decision support system}, CDSS) bertujuan
menutup celah tersebut dengan mencocokkan rekam medis pasien terhadap panduan klinis dan
menampilkan rekomendasi yang berlaku. Namun, CDSS yang efektif untuk CCD harus mengatasi
tiga tantangan yang masih ada. Pertama, sistem harus \emph{memadukan} tiga sumber bukti
yang berbeda---EHR pasien, panduan klinis, dan literatur biomedis pendukung---ke dalam satu
struktur yang dapat dikueri. Kedua, sistem harus \emph{menjembatani terminologi klinis yang
heterogen} (misalnya SNOMED~CT, RxNorm, dan LOINC) agar kode dalam rekam medis dapat
dicocokkan dengan konsep dalam panduan klinis. Ketiga, ketika disajikan dalam bahasa alami,
keluarannya harus \emph{tepercaya}, yaitu bebas dari klaim klinis yang dikarang ataupun
tidak berdasar pada data.

Terdapat dua perkembangan teknologi yang membuat pendekatan baru menjadi mungkin. Pertama, EHR modern semakin banyak memaparkan data klinis melalui standar HL7 \emph{Fast
Healthcare Interoperability Resources} (FHIR) \cite{fhir_r4}. FHIR merupakan format
terstruktur yang berorientasi pada sumber daya (\emph{structured, resource-oriented
format}). Format ini menjadikan kondisi, prosedur, obat, dan \emph{observasi numerik} (hasil
laboratorium serta tanda vital) sebagai data kelas satu yang dapat dibaca mesin. CDSS yang
bersifat \emph{FHIR-native} karenanya dapat mengonsumsi data klinis terstruktur secara
langsung, alih-alih bergantung pada ekspor tabular terdatarkan, yang di dalamnya observasi
berkomponen jamak seperti panel tekanan darah terurai menjadi baris-baris terpisah dan harus
dirakit ulang sebelum dapat dijadikan dasar penalaran (Subbab~\ref{sec:flattening}). Perakitan
ulang itu merupakan biaya, bukan penghalang; Subbab~\ref{sec:csv-gating} mengukur berapa
besarnya. Kedua, \emph{knowledge graph} (KG)---graf
pengetahuan yang diwujudkan di atas basis data \emph{property graph} seperti Neo4j---sangat
sesuai untuk menyatukan EHR, panduan klinis, dan literatur ke dalam satu graf, dengan keragaman
terminologi disatukan melalui pemetaan setiap kode ke kode payung bersama menggunakan
\emph{Unified Medical Language System} (UMLS) \cite{umls}.

Peluang sekaligus risiko yang saling melengkapi muncul dari \emph{large language model} (LLM). LLM mampu mengubah fakta rekomendasi yang
terstruktur menjadi ringkasan yang bermakna dan siap dibaca klinisi, tetapi rentan
terhadap \emph{hallucination}, yaitu pemunculan pernyataan yang tidak berdasar---sesuatu
yang tidak dapat diterima dalam konteks klinis. Rancangan yang penelusuran dan penalarannya bersifat \emph{deterministik}---dikerjakan
oleh kueri graf---dan yang peran LLM-nya dibatasi pada \emph{perumusan kalimat} atas
fakta hasil kueri beserta asal faktanya, dapat menghasilkan keterbacaan sekaligus
\emph{faithfulness} `kesesuaian keluaran terhadap fakta sumber' yang dapat diverifikasi. Hal ini berbeda dengan
\emph{retrieval-augmented generation} konvensional yang menelusuri potongan teks secara
aproksimatif melalui kemiripan vektor \cite{rag_lewis2020}.

Skripsi ini mengembangkan sistem rekomendasi klinis (\emph{clinical recommendation system}, CRS) berbasis \emph{knowledge graph}
yang bersifat FHIR-native untuk CCD. Rekam medis pasien sintetis namun realistis
dihasilkan dengan Synthea \cite{synthea} dan diekspor sebagai \emph{bundle} HL7 FHIR~R4,
kemudian dimuat ke dalam \emph{knowledge graph} Neo4j, lalu disatukan dengan panduan AHA/ACC CCD
2023 dan korpus literatur PubMed yang dihasilkan menggunakan \emph{seed} dari referensi panduan klinis tersebut melalui lapisan
konsep UMLS. Transformasi dari FHIR ke graf itu sendiri bukanlah hal yang baru dan telah tersedia
sebagai alat siap pakai yang berbasis web \cite{tenasol_viewer}; kontribusi penelitian ini terletak pada
lapisan-lapisan yang dibangun di atasnya. Kontribusi utamanya mengikuti langsung dari
rancangan FHIR-native. Karena observasi numerik dimuat sebagai data graf, sistem dapat
menerapkan \emph{observation gate} pada rekomendasi yang bergantung pada observasi
berdasarkan nilai terukur pasien yang sebenarnya. Contohnya adalah rekomendasi penyekat
beta yang disyaratkan pada fraksi ejeksi ventrikel kiri dan rekomendasi pengelolaan berat
badan yang disyaratkan pada indeks massa tubuh. Hal itu tidak dapat dilakukan oleh alur
berbasis kondisi dan teks semata. Perolehan tersebut dihasilkan oleh pemuatan nilai observasi numerik sebagai data kelas satu. FHIR merupakan cara yang paling langsung untuk memperolehnya, tetapi bukan satu-satunya cara. Subbab~\ref{sec:csv-gating} mengujinya dengan
membangun kembali gerbang yang sama dari ekspor tabular dan batas klaim yang sah menurut hasil
pengujian itu dinyatakan pada Subbab~\ref{sec:limitations}.
Sistem juga menghasilkan ringkasan dalam bahasa alami yang membumi (\emph{grounded briefing}),
yang isinya ditelusuri secara deterministik dari graf dan \emph{faithfulness}-nya pada tingkat sitasi diperiksa secara otomatis. Dengan demikian, penelitian ini menyelidiki bagaimana \emph{knowledge
graph} yang FHIR-native dapat meningkatkan presisi sekaligus tingkat kepercayaan
rekomendasi yang selaras dengan panduan klinis untuk penyakit jantung koroner kronis.

\section{Rumusan Masalah}

Berdasarkan latar belakang di atas, masalah yang hendak diatasi dalam penelitian ini
dirumuskan sebagai berikut:

\begin{enumerate}
    \item Bukti yang diperlukan untuk mendukung keputusan klinis pada penyakit jantung
    koroner kronis tersebar pada tiga sumber yang berbeda bentuk---rekam medis pasien,
    panduan praktik klinis, dan literatur biomedis---dan belum tersedia dalam satu
    struktur yang dapat dikueri. Keadaan tersebut diperberat oleh terminologi klinis yang
    heterogen: kondisi dikodekan dalam SNOMED~CT, obat dalam RxNorm, dan observasi dalam
    LOINC sehingga kode dalam rekam medis tidak dapat langsung dicocokkan dengan konsep
    yang dirujuk panduan klinis.

    \item Alur pendukung keputusan yang bertumpu pada kondisi dan teks semata tidak mengikutsertakan
    nilai observasi numerik, padahal sebagian rekomendasi dari panduan klinis disyaratkan pada nilai
    terukur pasien---misalnya fraksi ejeksi ventrikel kiri atau indeks massa tubuh.
    Akibatnya, rekomendasi yang bergantung pada observasi tidak dapat digerbangi menurut
    nilai observasi dari pasien yang sebenarnya dan keselarasan panduan klinis yang dihasilkan menjadi kurang
    presisi.

    \item Keluaran bahasa alami yang dihasilkan LLM rentan terhadap \emph{hallucination},
    sedangkan \emph{retrieval-augmented generation} konvensional hanya berlandas pada
    potongan teks yang ditelusuri secara aproksimatif. Belum tersedia rancangan yang
    \emph{faithfulness}-nya dapat diverifikasi terhadap rekam medis sumber---sifat yang
    diperlukan dalam konteks klinis.
\end{enumerate}

\section{Tujuan Penelitian}

Sejalan dengan rumusan masalah, tujuan penelitian ini adalah sebagai berikut.

\begin{enumerate}
    \item Merancang dan mengimplementasikan alur FHIR-native yang memuat \emph{bundle}
    HL7 FHIR~R4 ke dalam \emph{knowledge graph} Neo4j, lalu menyatukannya dengan panduan
    AHA/ACC CCD 2023 serta korpus literatur yang dibenihkan dari referensi pada panduan
    klinis melalui jembatan konsep berbasis UMLS. Sistem yang dihasilkan kemudian dievaluasi
    secara menyeluruh melalui uji keselarasan panduan klinis terhadap standar rujukan
    otomatis yang independen dari sistem rekomendasi, uji kesepakatan dengan standar emas
    berlabel manual (Cohen's~$\kappa$), deteksi keamanan kontraindikasi, serta telaah ketepatan
    rekomendasi oleh klinisi.
    \item Mengimplementasikan \emph{observation gate} yang mengevaluasi rekomendasi bergantung-observasi terhadap nilai
    terukur pasien. Perubahan keselarasan panduan klinis yang dihasilkan gerbang tersebut kemudian diukur melalui studi ablasi (dengan dan tanpa \emph{observation gate}).
    \item Mengimplementasikan alur \emph{grounded briefing} yang
    fakta-fakta rekomendasinya ditelusuri secara deterministik dari graf dan
    perumusan kalimatnya dilakukan oleh LLM. \emph{Faithfulness} ringkasan yang dihasilkan kemudian dievaluasi dengan prosedur verifikasi otomatis.
\end{enumerate}

\section{Batasan Penelitian}

Agar penelitian tetap terfokus, batasan berikut diberlakukan:

\begin{enumerate}
    \item \textbf{Ranah klinis.} Sistem hanya menyasar penyakit jantung koroner kronis
    dan mengodekan panduan AHA/ACC CCD 2023 \cite{virani2023} sebagai basis rekomendasinya.
    \item \textbf{Data.} Data pasien bersifat \emph{sintetis}, dihasilkan oleh Synthea
    dan diekspor sebagai \emph{bundle} HL7 FHIR~R4; tidak ada data pasien nyata yang
    digunakan. Kohort terdiri atas 5.490 pasien (lihat Bab~III) dan subhimpunan
    yang lolos dari gerbang CCD digunakan untuk evaluasi.
    \item \textbf{Terminologi.} Kondisi dan prosedur dikodekan dalam SNOMED~CT, obat dalam
    RxNorm, dan observasi dalam LOINC. Kode ICD-10-CM/ICD-11 tidak terdapat pada data
    sintetis; pemuatan rekam medis berkode ICD pada penerapan nyata akan memerlukan
    pemetaan silang SNOMED$\leftrightarrow$ICD tambahan.
    \item \textbf{Observasi.} Hanya observasi LOINC dalam \emph{allow-list} yang
    relevan secara klinis---misalnya fraksi ejeksi, LDL, indeks massa tubuh, trigliserida,
    tekanan darah, dan status merokok---yang dimuat dan digunakan oleh \emph{observation gate}.
    \item \textbf{Peran sistem.} Sistem merupakan alat \emph{rekomendasi} yang
    bersifat informatif dan menampilkan pilihan yang selaras dengan panduan klinis; sistem tidak
    mengambil keputusan klinis dan tidak diterapkan pada pelayanan pasien nyata.
    \item \textbf{Evaluasi.} Efektivitas diukur sebagai \emph{keselarasan panduan klinis}, bukan
    sebagai luaran klinis nyata. Sejumlah prasyarat pada panduan klinis yang tidak dapat dipastikan
    dari data terstruktur---misalnya skor risiko komposit dan sebagian klausa yang
    didefinisikan melalui negasi---diperlakukan sebagai keterbatasan yang terdokumentasi.
\end{enumerate}

\section{Manfaat Penelitian}

Manfaat yang diharapkan dari penelitian ini adalah sebagai berikut.

\begin{enumerate}
    \item \textbf{Manfaat akademis.}
    \begin{itemize}
        \item Menunjukkan arsitektur yang dapat direproduksi untuk CRS berbasis
        \emph{knowledge graph} yang FHIR-native dan menyatukan bukti dari EHR, panduan klinis,
        serta literatur.
        \item Mengukur secara kuantitatif kontribusi observasi terstruktur---yang
        dimungkinkan oleh EHR FHIR-native---terhadap presisi keselarasan panduan klinis.
        \item Memperkenalkan kerangka \emph{faithfulness} berbasis penelusuran deterministik untuk
        keluaran LLM pada konteks klinis, yang isi faktualnya dapat diverifikasi terhadap
        sumbernya.
    \end{itemize}
    \item \textbf{Manfaat praktis.}
    \begin{itemize}
        \item Menyediakan cetak biru untuk membangun sistem rekomendasi yang sesuai panduan klinis
        berbasis data FHIR terstandar, yang dapat dipindahkan ke kondisi dan panduan klinis lain.
        \item Menawarkan pola ringkasan membumi untuk pembangkitan teks klinis yang
        tepercaya, yang dapat menginformasikan---tanpa menggantikan---pertimbangan klinisi.
    \end{itemize}
\end{enumerate}

\section{Sistematika Penulisan}

Sisa skripsi ini disusun sebagai berikut.

\noindent\textbf{Bab~I (Pendahuluan)} memaparkan latar belakang, rumusan masalah, tujuan,
batasan, dan manfaat penelitian.

\noindent\textbf{Bab~II (Tinjauan Pustaka dan Dasar Teori)} menelaah penelitian terdahulu
mengenai pendukung keputusan klinis, \emph{knowledge graph} berbasis panduan klinis, dan
\emph{retrieval-augmented generation} serta menetapkan dasar teori yang mencakup HL7 FHIR,
basis data \emph{property graph}, pemetaan terminologi berbasis UMLS, dan konsep klinis
panduan CCD.

\noindent\textbf{Bab~III (Metode Penelitian)} menguraikan rancangan penelitian: sumber data
dan kohort, urutan pembangunan graf, rancangan pembentukan standar rujukan yang independen dari sistem
rekomendasi, serta protokol evaluasi.

\noindent\textbf{Bab~IV (Hasil dan Pembahasan)} melaporkan dan membahas kualitas
pembangunan graf, ablasi keselarasan \emph{observation gate}, validasi $\kappa$ terhadap
label manual, evaluasi deteksi keamanan, hasil \emph{faithfulness} ringkasan,
serta perbandingan terhadap penelitian terdahulu.

\noindent\textbf{Bab~V (Kesimpulan dan Saran)} merangkum temuan, menyatakan kesimpulan
terhadap tujuan penelitian, dan menguraikan arah penelitian selanjutnya.
